%\input ../util/bookmacros.tex
%\input ../util/docparams.tex

\chapter{Difference Calculus}

\section{Motivational Problem for Differencing}
Suppose there are several cars suspected in the use of an armed 
robbery, all of a particular make and model. Soon after the robbery,
a high-speed chase ensued between the police and the getaway car. The
car used had no license plates at the time of the crime; however, it's now believed that 
the getaway car has legitimate plates. There are not enough visual clues 
to distinguish the getaway car from the group of suspected cars. The police 
would like to be able to give enough evidence to a judge in order to 
get a search warrant for the owner of one of the cars. The one thing the 
police have going for them is that each car has a sensor which records the 
odometer reading every $1/10$ of a second. The odometer data is stored on an 
internal disk in feet. They want to examine 
the data in order to find out if there were high speeds occurring in any 
of these cars during the time of the getaway. 

Each car would have a 
very long list of numbers
representing odometer readings -- meaning distance. How should 
we analyze the data to find velocity? Well, from physics we know the 
average velocity over a time interval is the distance traveled divided by 
the time. So, the average velocity over each time interval of $1/10$th of 
a second is found by taking the difference in odometer readings over that 
interval and then dividing by $1/10$. Before discussing any more details, 
it would be convenient to have a notation to describe long lists of numbers 
in order to talk about operating on them. 

\section{Sequences}
We'll define a {\it sequence\/} of numbers to be a list of numbers. 
It is more than a collection of numbers; it is an ordered 
collection. The distinction is the same as when one considers an ordered list of tasks to perform, a ``to do'' list,
as opposed to an unordered collection of tasks.
For instance, the list $\{1,20,3,5,1,7\}$ is a sequence of numbers as is the list of
positive integers in increasing order; or the list of odd numbers listed in increasing order.
Since the sequence values 
are ordered we can talk about the $5^{\rm th}$ or the $10^{\rm th}$ element of a sequence (provided there are
that many elements). The $5^{\rm th}$ positive integer
is of course $5$; the $10^{\rm th}$ odd number is $19$.
In order to talk about sequences in more generality we use the notation 
$\{x_n\}_{n=1}^{N}$ to represent the sequence: 
$x_1, x_2, x_3, \ldots, x_N$; where the $x$ values are numbers and the three dots, $\ldots$, 
mean the intermediate values, in this case the sequence values between $x_3$ and $x_N$. 
Similarly, the notation $\{x_n\}_{n=0}^{N}$ is used 
to represent the sequence: $x_0, x_{1}, x_{2}, \ldots, x_N$. All that changed here 
is the choice of how to start off the indexing. More generally,
$\{x_n\}_{n=a}^{N}$ ($a$ is an integer) is used to represent the sequence:
$x_a, x_{a+1}, x_{a+2}, \ldots, x_N$. If there is no value listed after the dots then we mean that 
there is no end to the sequence -- like the sequence of odd numbers. 
In this case one would write: $\{x_n\}_{n=1}^{\infty}$.\footnote{\kern 0.5pt \raise 0.5ex \hbox{\dag}}%
{The symbol $\infty$ is meant to represent infinity.} The $5^{\rm th}$ and $10^{\rm th}$ values of 
such a sequence are $x_5$ and $x_{10}$ respectively.

\subsection{Sequences as Functions}
Sequences can be thought of as functions on the integers; they take an integer as input and return a number.
The functions that you are more familiar with, ``continuous'' functions like $f(x) = x^2$, take a number and return a number.
For both sequences and continuous functions there is a distinction to be made between the function and the formula -- if any --
that defines it. See the appendix ``Functions, Formulas, Graphs, and Notation'' for a discussion of this. 
An example of a sequence that can be described by a simple formula is the even numbers. This sequence is
$\{ 2,4,6,8,\ldots\}$. If we call this sequence $x$, then we can write the sequence values succinctly as $x_n = 2 n$.
In this case there is a formula that allows us to describe any element of the sequence simply. In cases where there is such 
a formula we will write the sequence using the formula. In the case of the even numbers we would write 
$\{2 n\}_{n=1}^{\infty}$. Consider now the sequence of primes: $\{ 2,3,5,7,11,13,17,\ldots\}$.
Calling this sequence $p$, we would like to write $p_n = $ some-nice-formula-in $n$. However, there is no such nice
formula. The sequence exists, but there is no known easily computable formula for it.


\section{Differencing of Sequences}
We know from the introduction that we need to ``difference'' the distance data 
to find the velocity. We define the difference symbol $\Delta$ 
(Greek letter Delta which corresponds to D -- for difference), 
to operate on sequences by differencing successive elements of a sequence.
Given a sequence, $\{x_n\}_{n=1}^{\infty}$, we define $\Delta x_n$ by:
$$
\Delta x_n = x_{n+1} - x_n \quad {\rm for} \; n \ge 1
$$
When using $\Delta$ on a sequence, $x$,  we may replace $x_n$ by its formula -- if it exists.
For instance, if $x_n = 5n+1$, then we might write $\Delta x_n$ as $\Delta (5n+1)$.
For this sequence, $\Delta x_n = x_{n+1} - x_n = (5(n+1) + 1) - (5n+1) = 5$.
If $p$ is the sequence of primes then what is $\Delta p_n$? Well, there is no formula for $p_n$ so we 
can't say -- meaning we can't write down a formula for this. However, given a particular $n$, say $7$, $\Delta p_7 = p_8 -
p_7$. Since the $8^{\rm th}$ prime is $19$ and the $7^{\rm th}$ prime is 17, we have 
$\Delta p_7 = 19 - 17 = 2$.

We can think of $\Delta$ as taking a \index{sequence} as input and returning a new sequence as the output.
If we are given a finite sequence, then one must be careful at 
the end point. For a finite sequence, $\{x_1, x_2, \ldots, x_N\}$, 
$\{x_n\}_{n=1}^{N}$, in our notation, $\Delta$ is defined to act 
on this sequence by
$$
\Delta x_n = x_{n+1} - x_n \quad {\rm for} \; n \in [1, N-1]
$$
In other words, the differencing is only defined up to $N-1$. That is,
up to where the differencing makes sense.

\example{If our sequence is $\{1,3,7,15,22\}$, then applying 
the differencing operation $\Delta$ to this sequence yields a new sequence: 
$\{2,4,8,7\}$.}

Notice that this sequence has one fewer element than the original.

\example{If $\{x_n\}_{n=1}^\infty$ is the list of odd numbers 
listed in increasing order, then after applying $\Delta$ to the sequence 
we obtain a new sequence: $\{2, 2, 2, 2, 2, \ldots\}$.}

\section{Sample Difference Formulas}
We will soon want a stock of difference formulas on hand; in the next chapter
we will be running $\Delta$ in reverse, and the more differences we recognize
on sight, the easier that game becomes.
Here are a few difference formulas for sequences of successive powers, along with
the difference of a constant sequence.
Let $w_n = c$ (a constant sequence), $x_n = n$, $y_n = n^2$, and $z_n  = n^3$, then
$$
\eqalign{
\Delta w_n = \Delta c & = c - c = 0 \cr
\Delta x_n = \Delta n & = (n+1) - n  = 1 \cr
\Delta y_n = \Delta n^2 & = (n+1)^2 - n^2 = n^2 + 2 n + 1 - n^2 = 2 n + 1 \cr
\Delta z_n = \Delta n^3 & = (n+1)^3 - n^3 = n^3 + 3 n^2 + 3 n + 1 - n^3 = 3 n^2 + 3 n + 1
}
$$
Notice that the right hand side of the above formulas is a polynomial of one less power than the
left hand side.

\exercise{Show that $\Delta 2^n = 2^n$; that is, the sequence $\{2^n\}_{n=0}^\infty$
is its own difference. (Much later, in the chapter on Differential Calculus, we
will meet a famous function with the analogous property.)}

\exercise{Compute the difference sequence of $\{0,\,1,\,4,\,9,\,16,\,25\}$. Then
compute the difference sequence of {\it that\/} sequence -- the ``second
difference''. What do you notice?}

\section{Back to finding the maximum velocity}
We can express the average velocity (in feet per second) of the car over the $n^{\rm th}$ interval
$[n/10, (n+1)/10]$ as: $v_n = \Delta x_n / (1/10) = 10 \, \Delta x_n$. So, we have a new
sequence, $\{v_n\}$ which represents an approximation to the velocity.
More generally, if $h$ is the length of the time interval, then the sequence 
of approximate velocities may be computed as:
$$
v_n = \Delta x_n / h
$$
To find the maximum velocity we could write a program which walks the sequence,
$\{v_n\}$, looking for the maximum. At each point one compares the maximum 
value found so far with the current data point. If the data point has a larger value, 
we change the maximum value to this value. Such a program might take an array 
(representing the sequence of velocities) and return the maximum value. 
We offer some code written in the programming language equivalent of 
	\index{Esperanto}%
\footnote{\kern 0.5pt \raise 0.5ex \hbox{\ddag}}{We use the characters // to denote 
the beginning of a comment line.}.
\vskip 12pt
{\tt \obeylines \obeywhitespace \parskip=0pt
find-max(v) $\lbrace$
   // The current maximum - so far.
   vmax := v[0]

   // The number of values to examine.
   N := length(v)

   // Loop over the rest of the values and compare...
  {\bf for\/} i {\bf in\/} [1,N-1] {\bf do\/}

      // Update current maximum if bigger value found.
      {\bf if\/} (vmax $<$ v[i]) {\bf then\/}
         vmax := v[i] 
      {\bf end if\/}
   {\bf end for\/}
  
   // Return the maximum.
  {\bf return\/} vmax 
$\rbrace$
}

Notice the work involved in finding the maximum \index{velocity}: the loop examines
all $N$ values, making $N-1$ comparisons along the way.

After running this program on the data from the cars the police crime lab is
able to identify
a car which hit speeds of $80$ mph. (Since the odometer data is in feet and
the time steps are in seconds, the program reports velocities in feet per
second; the lab converts to miles per hour by multiplying by ${3600 \over 5280}$
-- so $80$ mph is about $117$ ft/sec.) Unfortunately, the manufacturer informs them
that while the dash-board odometer reading the owner sees is correct, 
the {\it historical\/} odometer data stored on disk is somewhat inaccurate.
This ``historical'' odometer data was designed as a tracking system,
originally intended to defeat attempts by used car dealers to illegally
modify odometer readings, and it promised to help the service departments
maintain the cars as well. But the manufacturer worried about privacy:
no one wants their every move recorded, and lawsuits -- or simply bad
press -- could follow. The engineering department found an ingenious
compromise. They would alter the stored data so that it was a bit fuzzy --
still useful to the service departments, who needed only general driving
patterns along with velocity and acceleration over short intervals. There
was a secondary benefit: a small signal could be embedded in the modified
data to act as a fingerprint, making it easier to spot if someone had
tampered with the recorded data.

This left the police with a problem. They asked the auto manufacturer for
some information about its process of ``fuzzying up'' the data; after all,
all the police wanted to know was whether one of the cars went over $80$
mph on the day of the robbery. The auto maker said the details were a
Trade Secret, but agreed to say roughly what the process was. At time step $n$ the 
odometer value, $x_n$, is modified by multiplying it by a factor which depends 
on $n$. Let us call this factor $f_n$, then the value stored on 
the disk is $d_n = x_n f_n$. Although the auto maker would not 
tell what $f_n$ 
was, they would say that it varied in a range from  
$0.5$ to $2.0$ back to $0.5$. It varied in this way over a period greater 
than $10,000,000$ time steps. This meant the period was over 10 days 
of continuous driving. 

What the police have recorded on disk was not the odometer reading, $x_n$, 
but $d_n = x_n f_n$. They needed a way to relate the differencing of 
$x_n f_n$ to the differencing of $x_n$ (Then it is a simple matter to 
get the velocity -- divide by $h$). 

The difference of the disk data is $\Delta d_n = \Delta (x_n f_n) = x_{n+1} f_{n+1} - x_n f_n$.
We would like to relate this difference to the difference of the actual distance values: $\Delta x_n$.
This is accomplished by adding and subtracting the term $x_n f_{n+1}$.
$$
\eqalign{
\Delta (x_n f_n) & = x_{n+1} f_{n+1} - x_n f_n \cr
& = x_{n+1} f_{n+1} - x_n f_{n+1} 
+ x_n f_{n+1} - x_n f_n \cr
& = (x_{n+1} - x_n) f_{n+1} + x_n (f_{n+1} - f_n) \cr
& = \Delta x_n \, f_{n+1} + x_n \Delta f_n \cr
}
$$
Or,
$$
\eqalign{
\Delta (x_n f_n) & = \Delta x_n f_{n+1} + x_n \Delta f_n \cr
}
$$
This gives us the relationship we were after. 
Since $x_n f_n  = f_n x_n$ we also have
$$
\eqalign{
\Delta (x_n f_n)  &= \Delta (f_n x_n)  = \Delta f_n x_{n+1} + f_n \Delta x_n \cr
}
$$
Let us formalize what we have so far as a theorem. We state the result with different
names for the sequences.
\proclaim Theorem(Discrete Product Rule). If $\{x_n\}_{n=0}^{\infty}$ and 
$\{y_n\}_{n=0}^{\infty}$ are sequences then the difference of 
their product is related to the individual differences by:
$$
\eqalign{
\Delta (x_n y_n) & = \Delta x_n y_{n+1} + x_n \Delta y_n \cr
}$$
Or,
$$
\eqalign{
\Delta (x_n y_n) & = x_{n+1} \Delta y_{n} + \Delta x_n y_n 
}
$$

Getting back to our problem, we wish to find $\Delta x_n$ from $\Delta (x_n f_n)$ (the difference of
the recorded odometer values),
and we have the relation: $\Delta (f_n x_n)  = \Delta (x_n f_n) = \Delta x_n f_{n+1} + x_n \Delta f_n$.
Ideally, we would like to solve this equation for
$\Delta x_n$.
This is the same game as solving $3x + 2 = 8$ in algebra: get the unknown by itself
under one operation, then undo that operation -- a move we'll call {\it isolate and invert\/}.
Something like:
$$
\Delta x_n f_{n+1} = \Delta (x_n f_n) - x_n \Delta f_n
$$
So that
$$
\Delta x_n  = \bigl( \Delta (x_n f_n) - x_n \Delta f_n \bigr) / f_{n+1}
$$
Unfortunately, the formula for $\Delta x_n$ involves $x_n$, $f_n$, and $\Delta f_n$ which are all unknown.
But we are not necessarily looking for an exact solution; we would be 
happy with a statement which said $v_n = \Delta x_n / h \ge 80$ for some value of $n$. This 
is actually more incriminating than saying $v_n = 80$. 

Since $f$ changes over a very long scale in comparison with 
$x$, it is essentially constant for the purposes of our calculations. That is, 
for all intents and purposes over the short time span of the robbery $f_n$ can be treated as a constant
equal to some value, say $F$. Replacing a slowly varying quantity by a constant is an
approximation -- the first half of the {\it approximate, then take a limit\/} move
we will use throughout this book. If a sequence, $f_n$, is constant
then $\Delta f_n = 0$. Therefore, the last equation becomes%
\footnote{\kern 0.5pt \raise 0.5ex \hbox{\dag}}{The symbol $\approx$
  means approximately equal.}
$$
\Delta x_n  \approx \Delta (x_n f_n) / F
$$
We still don't know $F$. Here it is a simple matter for the police to 
run the car for a short distance and compare the odometer tracking data with
the actual distance traveled. Because $f_n$ has been virtually constant since the robbery (a few hours before)
they are able to run an experiment to find $F$. They find that it 
was approximately $0.8$ for the suspect car during the robbery. 
This means their maximum velocity calculation must be corrected by dividing by 0.8 for 
the suspect car. 
In this case $v_n = \Delta x_n /h \approx \Delta(x_n f_n) / (F h) \approx  80 / 0.8$ mph = $100$ mph. 
The police now have enough 
evidence to get a search warrant for the premises of the owner of this car.

\exercise{Find a similar formula expressing the difference of the
product $x_n y_n z_n$ in terms of the differences:
$\Delta x_n$, $\Delta y_n$, and $\Delta z_n$.}

\section{Summary}
We started with a list of odometer readings and a question about velocity, and
ended up inventing the difference operation. Along the way we made the idea of a
list precise with sequences and the notation $\{x_n\}_{n=1}^{N}$, and we saw that
differencing a sequence produces a new sequence -- one that measures how fast the
original changes.

Here is what is now in our toolkit:
\beginEnum
\enum{The definition of the difference operation: $\Delta x_n = x_{n+1} - x_n$.}
\enum{Sample difference formulas: $\Delta c = 0$, $\Delta n = 1$, $\Delta n^2 = 2n + 1$,
and $\Delta n^3 = 3n^2 + 3n + 1$.}
\enum{The Discrete Product Rule: $\Delta (x_n y_n) = \Delta x_n \, y_{n+1} + x_n \Delta y_n$.}
\enum{The {\it isolate and invert\/} strategy: to recover an unknown from an equation,
get it by itself under one operation and then undo that operation.}
\endEnum
In the next chapter we turn the difference operation around and ask the reverse
question: given the differences, how do we recover the original sequence?




